% Source-derived chapter 04: Voisin's Theorem in Odd Genus
% Original source: universal-secant-bundles-syzygies-canonical-curves
\chapter{Voisin's Theorem in Odd Genus}
\label{universal-secant-bundles-syzygies-canonical-curves-chapter-04}
\source{universal-secant-bundles-syzygies-canonical-curves}

\begin{remark}[Source section]
\label{universal-secant-bundles-syzygies-canonical-curves-chapter-04-source}
\source{universal-secant-bundles-syzygies-canonical-curves}
This chapter records the source section; no Lean declaration has yet been checked for this material.
\notready
\end{remark}

% The source section title was: Voisin's Theorem in Odd Genus
\label{odd-genus}
\source{universal-secant-bundles-syzygies-canonical-curves}
In this section we prove Green's Conjecture for generic curves of odd genus $g=2k+1 \geq 9$. We broadly follow a similar strategy as \cite{V2} to reduce to the even genus case, however the Secant Bundle approach provides significant simplifications. As in Voisin's proof, we work with a K3 surface $X$ of Picard rank two over $\C$, with $\text{Pic}(X)$ generated by a big and nef line bundle $L'$ with $(L')^2=2k+2$ together with a smooth rational curve $\Delta$ with $(L' \cdot \Delta)=0$.

The following lemma is adapted from \cite[Prop.\ 1]{V2}.
\begin{lem}
\source{universal-secant-bundles-syzygies-canonical-curves}
\notready \label{BNP-gen}
 Set $L_n:=L'-n\Delta$ for $0 \leq n \leq 3$.  Then $L_n$ is base-point free and is further ample for $0<n \leq 2$. Furthermore, the linear system $|L_n|$ cannot be written as a sum of two pencils. In particular, smooth curves $C \in |L_n|$ are Brill--Noether general. If $C \in |L_n|$ is general then $C$ is, in addition, Petri general.
\end{lem}
\begin{proof}
%If $\Delta$ is not integral, we have a smooth rational component $R$ of $\Delta$ with $(R \cdot \Delta)<0$. Writing $R=aL'+b\Delta$ then, since $L'$ is nef, and $R, \Delta-R$ are effective, $a=0$. But $(R)^2=-2$ so $b=1$ and $R=\Delta$, which is a contradiction.
Since $(L_n)^2 \geq 0$ and $(L_n \cdot L')>0$, we see $\dim |L_n|>0$ by Riemann--Roch. We claim $L_n$ is nef, i.e.\ there is no rational, base component $R \sim aL'+b\Delta $ of $L_n$ with $(L_n \cdot R)<0$. Otherwise, $R$ and $L_n-R$ would be effective and $a \neq 0$ so $a =1$, but then $\dim |L_n|=\dim|L_n-R|=0$. Since $(\alpha \cdot \beta)$ is even for all $\alpha, \beta \in \mathrm{Pic}(X)$, there are no divisors $E$ with $(L_n \cdot E)=1$, hence the nef bundles $L_n$ are base-point free, \cite[Prop.\ 8]{mayer}. If $L_n$ is not ample for $n \neq 0$, there is a rational curve $R\sim aL'+b\Delta$ with $(L_n \cdot R)=0$. Then $b=\frac{-ag}{n}$ with $a>0$, so $(R)^2=-2a^2g(\frac{g}{n^2}-1)<-2$.

For the last claim it suffices to show that we cannot write $L_n=A_1+A_2$ for divisors $A_i$ with $h^0(A_i)\geq 2$, $i=1,2$, \cite{lazarsfeld-BNP}. Writing $A_i=a_i L'+b_i \Delta$, we must have $a_j=0$ for some $j\in \{1,2 \}$. But then $h^0(A_j)=1$ which is a contradiction.
\end{proof}

Setting $L:=L_1=L'-\Delta$, a general $C \in |L|$ is a curve of genus $g=2k+1$. To verify Green's conjecture for $C$, it suffices to show $\mathrm{H}^1(\wedge^{k+1}M_L)=\mathrm{K}_{k,1}(X,L)=0$. We have an exact sequence $0 \to M_L \to M_{L'} \to \mathcal{O}(-\Delta) \to 0$
of vector bundles on $X$. This gives an exact sequence $$0 \to \wedge^{k+1}M_L \to \wedge^{k+1}M_{L'} \to \wedge^k M_L (-\Delta) \to 0.$$
The induced map $pr_{k}: \; \mathrm{H}^1(\wedge^{k+1}M_{L'}) \to \mathrm{H}^1(\wedge^k M_L (-\Delta))$ is called the \emph{projection map}.
\begin{lem}
\source{universal-secant-bundles-syzygies-canonical-curves}
\notready \label{proj-crit}
Suppose the projection map $pr_k$ is injective. Then $\mathrm{H}^1(\wedge^{k+1}M_L)=0$.
\end{lem}\begin{proof}
Indeed $\mathrm{H}^0(\wedge^k M_L (-\Delta))\seq \mathrm{H}^0(\wedge^k M_L)=0$. %From the defining sequence $0 \to M_L \to \mathrm{H}^0(L) \otimes \mathcal{O}_X \to L \to 0$, we obtain a long exact sequence
%$$0 \to \wedge^k M_L \to \wedge^k \mathrm{H}^0(L) \otimes \mathcal{O}_X \to \wedge^{k-1}\mathrm{H}^0(L) \otimes L \to \wedge^{k-2}\mathrm{H}^0(L) \otimes L^2 \to \ldots. $$
%The claim now follows from the fact that $\delta : \wedge^k \mathrm{H}^0(L) \to \wedge^{k-1}\mathrm{H}^0(L) \otimes \mathrm{H}^0(L)$ is injective, since if $\wedge : \wedge^{k-1}\mathrm{H}^0(L) \otimes \mathrm{H}^0(L) \to \wedge^k \mathrm{H}^0(L)$ is the natural multiplication map, we have $\wedge \circ \delta = \pm k \mathrm{Id}$.
\end{proof}
%We will prove that $pr_{k}: \; \mathrm{H}^1(\wedge^{k+1}M_{L'}) \to \mathrm{H}^1(\wedge^k M_L (-\Delta))$ is injective.

%The image under the morphism $\phi_{L'}: X \to \PP^{g+1}$ is a nodal K3 surface $\hat{X}$.
 We have a rational resolution of singularities $\mu : X \to \hat{X}$, contracting $\Delta$ to a du Val singularity $p$ on a nodal K3 surface $\hat{X}$. Then $\hat{X}$ admits a line bundle $\hat{L}$ with $\mu^*{\hat{L}}=L'$. We have $\mathrm{Pic}(\hat{X})=\mathrm{Cl}(\hat{X})\simeq \mathrm{Cl}(X\setminus \Delta)$, so  $\mathrm{Cl}(\hat{X})=\mathbb{Z}[\hat{L}]$ and $\hat{X}$ is factorial. Consider the rank two Lazarsfeld--Mukai bundle $\hat{E}$ on $\hat{X}$ induced by a $g^1_{k+2}$ on a general $C \in |\hat{L}|$. Set $E:=\mu^*\hat{E}$, which is a rank two bundle on $X$ induced by a $g^1_{k+2}$ on a general $C \in |L'|$.

\begin{lem}
\source{universal-secant-bundles-syzygies-canonical-curves}
\notready \label{green-pic2}
We have $\mathrm{K}_{k+1,1}(X,L')=0$. Further, there is a natural isomorphism $$\mathrm{Sym}^{k-1}\mathrm{H}^0(X,E) \simeq \mathrm{K}_{k,1}(X,L').$$
\end{lem}
\begin{proof}
Set $\PP=\PP(\mathrm{H}^0(\hat{X}, \hat{E}))$ and let $\hat{\mathcal{Z}} \seq \hat{X} \times \PP$ be the locus defined by $\{ (x,s) \; | \; s(x)=0 \}$. Then $\hat{\mathcal{Z}}$ is a projective bundle over $\hat{X}$, is a local complete intersection in $\hat{X} \times \PP$ and is finite and flat over $\PP$ by ``miracle flatness'', \cite[Prop.\ 6.1.5]{EGA}. Note that, as in Section \ref{sec1}, finiteness of $\hat{\mathcal{Z}}$ over $\PP$ comes from the fact that $\mathrm{Cl}(\hat{X})=\mathbb{Z}[\hat{L}]$ and $h^0(\hat{E}(-\hat{L}))=h^0(\hat{E}^{\vee})=0$. The formal cohomological arguments from Section \ref{sec1} and Section \ref{sec2}, up to Corollary \ref{main-cor-nat-iso}, go through unchanged. Hence we see $\mathrm{K}_{k+1,1}(\hat{X}, \hat{L})=0$, and $\mathrm{Sym}^{k-1}\mathrm{H}^0(\hat{E}) \simeq \mathrm{K}_{k,1}(\hat{X},\hat{L})$. Since $\hat{X}$ has rational singularities, $\mathrm{H}^0(X,nL') \simeq \mathrm{H}^0(\hat{X}, n\hat{L})$ for all $n$, and the claim follows.
\end{proof}

We will prove that we have a natural injection $\mathrm{Sym}^{k-1}\mathrm{H}^0(X,E) \hookrightarrow \mathrm{H}^1(\wedge^k M_L (-\Delta))$. By the above lemma, this implies injectivity of $pr_k$, so that Voisin's Theorem follows.
\begin{lem}[Voisin, \cite{V2}, Proposition 2]
\source{universal-secant-bundles-syzygies-canonical-curves}
\notready \label{decomposable}
The natural map $d: \wedge^2 \mathrm{H}^0(E) \to \mathrm{H}^0(\wedge^2 E)=\mathrm{H}^0(L')$ does not vanish on decomposable elements.
\end{lem}
\begin{proof}
Suppose $v_1, v_2 \in \mathrm{H}^0(E)$ with $v_1 \wedge v_2 \neq 0$, $d(v_1 \wedge v_2)=0$. Then $v_1, v_2$ generate a rank one subsheaf $H_1 \seq E$ with at least two sections. Let $H_2:=E/H_1$ and $M_2:=H_2/\mathrm{T}(H_2)$, where $\mathrm{T}(H_2)$ denotes the torsion subsheaf. We have a short exact sequence $0 \to M_1 \to E \to M_2 \to 0$, with $M_1,M_2$ rank-one, torsion-free sheaves and $h^0(M_1) \geq 2$. Then $M_i=N_i \otimes I_i$, for $N_i$ a line bundle and $I_i$ an ideal sheaf of a $0$-dimensional scheme for $i=1,2$. Since $E$ is globally generated with $h^2(E)=0$, we have $h^0(N_2) \geq 2$. But then $\det(E)=L'=N_1 \otimes N_2$ is a sum of line bundles with at least two sections. We have already seen that this cannot occur.
\end{proof}
To set things up, we need a lemma.
\begin{lem}
\source{universal-secant-bundles-syzygies-canonical-curves}
\notready
The bundle $E(-\Delta)$ is isomorphic to the Lazarsfeld--Mukai bundle $F$ corresponding to a $g^1_{k}$ on a general $C' \in |L-\Delta|$. In particular, $E(-\Delta)$ is globally generated.
\end{lem}
\begin{proof}
We claim that $F$ is $\mu$-stable with respect to $L-\Delta$. Otherwise, we have a filtration
$$0 \to M \to F \to N \otimes I_{\zeta} \to 0$$
where $M, N$ are line bundles, $I_{\zeta}$ is the ideal sheaf of a $0$-dimensional scheme of length $k-(M \cdot N)$, where $h^0(N) \geq 2$ and with $\mu(M) \geq \mu(F)=g-4 \geq \mu(N)$, cf.\ \cite{lelli-chiesa}. We have $h^2(M)=0$ since $\mu(M)>0$ and further $(M)^2=\mu(M)-(M\cdot N)\geq (g-4)-k \geq 0$. Thus $h^0(M) \geq 2$ by Riemann--Roch, contradicting that $L-\Delta=\det(F)$ cannot be written as a sum of two pencils.


The rank two bundle $E(-\Delta)$ is simple with $\det E(-\Delta)=L-\Delta$ and $\chi(E(-\Delta))=k+1$. We claim that $E(-\Delta)$ is stable with respect to $L-\Delta$. Since $F$ is the unique stable bundle with these invariants, $E(-\Delta) \simeq F$. If $E(-\Delta)$ is not stable, choose a filtration $0 \to M' \to E(-\Delta) \to N' \otimes I_{\zeta'} \to 0$
as above. We again have $h^0(M') \geq 2$, and since $N'(\Delta) \otimes I_{\zeta'}$ is a quotient of the globally generated bundle $E$ and $h^2(E)\neq 0$, $h^0(N'(\Delta)) \geq 2$. This contradicts that $L$ cannot be written as a sum of two pencils.
\end{proof}



Following \cite[pg.\ 1177]{V2}, let $N$ denote the rank $k-1$ kernel bundle fitting into the sequence
$$0 \to N \to \mathrm{H}^0(E(-\Delta))\otimes \mathcal{O}_X \to E(-\Delta) \to 0.$$
For any $t \in \mathrm{H}^0(E)$, we have a map $\wedge \, t :\mathrm{H}^0(E(-\Delta)) \to \mathrm{H}^0(\wedge^2 E (-\Delta))=\mathrm{H}^0(L)$, inducing a map $N \to M_L$. This globalizes to a vector bundle morphism
$$r: \; N  \boxtimes \mathcal{O}_{\PP}(-1) \hookrightarrow \mathcal{M} ,$$
on $X \times \PP$, where $\PP:=\PP(\mathrm{H}^0(X,E)) \simeq \PP^{k+2}$, where $p: X \times \PP \to X$, $q: X \times \PP \to \PP$ are the projections and $\mathcal{M}:=p^*M_L$.
\begin{remark}
\source{universal-secant-bundles-syzygies-canonical-curves}
\notready
We may compare the definition of $N$ to the vector bundles defined in Section \ref{sec1}. Let $(X,L)$ be a polarized K3 surface of even genus $g=2k$. Then, for any $t \in \mathrm{H}^0(X,E)$ we have an isomorphism between the secant sheaf $\mathcal{S}'_t$ and
$$N_E:= \mathrm{Ker}(\mathrm{H}^0(E) \otimes \mathcal{O}_X \twoheadrightarrow E),$$
as follows readily from the exact sequence $0 \to \mathcal{O}_X \xrightarrow{t} E \xrightarrow{\wedge t} I_{Z(t)} \otimes L \to 0$. Notice that the isomorphism $\mathcal{S}'_t \simeq N_E$ depends on $t$.
\end{remark}

From Lemma \ref{decomposable}, $r$ fails to be injective \emph{on fibres} precisely at points in the locus
$$\mathcal{Z}:=\left \{ (x,s) \in X \times \PP\left(\mathrm{H}^0(E(-\Delta))\right) \; | \; s(x)=0 \right\}.$$
Thus $\mathrm{Coker}(r)$ fails to be locally free along $\mathcal{Z} \seq X \times \Lambda$, where $\Lambda:=\PP\left(\mathrm{H}^0(E(-\Delta))\right) \seq \PP$. \\

We rectify the failure of $\mathrm{Coker}(r)$ to be locally free through a standard construction. Define $\pi: B \to X \times \PP$ as the blow-up along the codimension four locus $\mathcal{Z}$ and let $p':=p\circ \pi$, $q':=q \circ \pi$. Let $j: D \hookrightarrow B$ be the exceptional divisor. For any vector bundle $A$ on $X \times \PP$,  \begin{align} \label{mult-D-iso}
\source{universal-secant-bundles-syzygies-canonical-curves}
\mathrm{H}^j(B,\pi^*A (nD))\simeq \mathrm{H}^j(B,\pi^*A )\end{align} for any $j$ and for $1 \leq n \leq 3$.


For $p=(x,t) \in \mathcal{Z}$, the kernel of $r \otimes k(p)$ is isomorphic to $\C\langle t\rangle \seq \mathrm{H}^0(X,E)$. Thus $$\mathrm{Ker}(r_{|_{\mathcal{Z}}})\simeq ({q}^*\mathcal{O}_{\PP}(-2))_{|_{\mathcal{Z}}},$$ where the inclusion $({q}^*\mathcal{O}_{\PP}(-2))_{|_{\mathcal{Z}}} \hookrightarrow (N  \boxtimes \mathcal{O}_{\PP}(-1))_{|_{\mathcal{Z}}}$ is given by the section $u \in \mathrm{H}^0(\mathcal{Z},N  \boxtimes \mathcal{O}_{\PP}(1))\seq \mathrm{H}^0(E) \otimes \mathrm{H}^0({q}^*\mathcal{O}_{\PP}(1)_{|_{\mathcal{Z}}})$ obtained by restricting $id \in \mathrm{H}^0(E) \otimes \mathrm{H}^0({q}^*\mathcal{O}_{\PP}(1))$ to $\mathcal{Z}$.

We now perform an \emph{elementary transformation} on  $N  \boxtimes \mathcal{O}_{\PP}(-1)$. Define $S$ as the dual bundle to $\mathrm{Im}({\pi^*r}^{\vee})$. Then $S^{\vee}$ is a vector bundle of rank $k-1$ fitting into the exact sequence
$$0 \to S^{\vee} \to \pi^* (N^{\vee} \boxtimes \mathcal{O}_{\PP}(1)) \to {{q'}^* \mathcal{O}_{\PP}(2)}_{|_D} \to 0.$$
From the definition of $S$, we have an exact sequence $$0 \to S \to \pi^* \mathcal{M} \to \Gamma \to 0,$$ where $\Gamma$ is locally free of rank $k+2$.

The blow-up of a projective space $\PP(V)$ along a subspace $W \seq V$ is a projective bundle over $\PP(V/W)$, \cite[\S 9.3.2]{3264}. Thus $B$ is a projective bundle $\PP(\mathcal{H})$ over $\PP(\mathcal{F})$, where $\mathcal{F}:=\mathrm{Coker}(N \hookrightarrow \mathrm{H}^0(E) \otimes \mathcal{O}_X)$. %Thus $B$ is the closure of the graph of the relative projection $\PP(\mathrm{H}^0(E)\otimes \mathcal{O}_X) \dashrightarrow \PP(\mathcal{F})$ away from $N$, \cite[Prop.\ IV-22]{eisenbud-harris}.
The projection morphism ${\chi: B \to \PP(\mathcal{F})}$ is defined over $X$ with $$\chi^* \mathcal{O}_{\PP(\mathcal{F})}(1) \simeq {q'}^*\mathcal{O}_{\PP}(1)(-D).$$ To describe $\mathcal{H}$, let $f: \PP(\mathcal{F}) \to X$ be the projection and define $\mathcal{P}$ by the exact sequence $0 \to \mathcal{P}^{\vee} \to f^*f_* \mathcal{O}_{\PP(\mathcal{F})}(1) \to \mathcal{O}_{\PP(\mathcal{F})}(1) \to 0.$ We have a surjection $\mathrm{H}^0(E) \otimes \mathcal{O}_{\PP(\mathcal{F})} \to \mathcal{P}$. The rank $k$ bundle $\mathcal{H}$ is defined by the exact sequence
$$0 \to \mathcal{H} \to \mathrm{H}^0(E) \otimes \mathcal{O}_{\PP(\mathcal{F})}\to \mathcal{P} \to 0.$$
We have a short exact sequence
$$ 0 \to f^* N \to \mathcal{H} \to \mathcal{O}_{\PP(\mathcal{F})}(-1) \to 0.$$
The isomorphism $B \simeq \PP(\mathcal{H})$ gives an identification $\mathcal{O}_{\PP(\mathcal{H})}(1) \simeq {q'}^*\mathcal{O}_{\PP}(1)$, \cite[Prop.\ 9.11]{3264}.
In summary, we have the following commutative diagram on $\PP(\mathcal{F})$:
{\small{$$\text{[Diagram omitted; see source PDF.]},$$}}
Note that we are using the geometric notation for projective bundles, i.e.\ $\PP(\mathcal{F})=\mathrm{Proj}(\mathcal{F}^{\vee})$, and $\mathcal{F}^{\vee}=f_*\mathcal{O}_{\PP(\mathcal{F})}(1)$.
\begin{lem}
\source{universal-secant-bundles-syzygies-canonical-curves}
\notready
We have  $S \simeq T_{\chi} \otimes {q'}^*\mathcal{O}_{\PP}(-2)$, where $T_{\chi}$ is the relative tangent bundle.
\end{lem}
\begin{proof}
We keep track of the maps defined above in the commutative diagram
{\small{$$\text{[Diagram omitted; see source PDF.]}.$$}}
We have the relative Euler sequence $0 \to \mathcal{O}_{\PP(\mathcal{H})} \to \mathcal{O}_{\PP(\mathcal{H})}(1) \otimes \chi^* \mathcal{H} \xrightarrow{\alpha} T_{\chi} \to 0.$
Twisting by $\mathcal{O}_{\PP(\mathcal{H})}(-2)$, we have the composite map
$\pi^*(N \boxtimes \mathcal{O}_{\PP}(-1)) \to \chi^*\mathcal{H} \otimes {q'}^*\mathcal{O}_{\PP}(-1)\xrightarrow{\alpha} T_{\chi}(-2).$
Dualizing, we obtain an exact sequence
$$0 \to \Omega_{\chi}\otimes {q'}^*\mathcal{O}_{\PP}(2) \to \pi^*(N^{\vee} \boxtimes \mathcal{O}_{\PP}(1)) \to {{q'}^*\mathcal{O}_{\PP}(2)}_{|_D} \to 0,$$
and comparison with the defining sequence for $S^{\vee}$ gives $S^{\vee} \simeq  \Omega_{\chi}\otimes {q'}^*\mathcal{O}_{\PP}(2)$.
%We have $\mathcal{Z}=\PP(N)$, and from the short exact sequence $0 \to N \to \mathrm{H}^0(E) \otimes \mathcal{O}_X \to \mathcal{F} \to 0$, we see that $I_{\mathcal{Z}/X \times \PP}$ is the image of a natural map $\mathcal{F} \boxtimes \mathcal{O}_{\PP}(-1) \to \mathcal{O}_{X \times \PP}$. So we have a morphism ${q'}^*\mathcal{O}_{\PP}(-1) \to {p'}^* \mathcal{F}$ which vanishes precisely on $D$, providing a surjection ${p'}^* \mathcal{F}^{\vee} \to {q'}^*\mathcal{O}_{\PP}(1)(-D)$. This gives us a morphism $\chi': B \to \PP(\mathcal{F})$ over $X$ with ${\chi'}^*\mathcal{O}_{\PP(\mathcal{F})}(1)\simeq {q'}^*\mathcal{O}_{\PP}(1)(-D)$.
\end{proof}

We will need the following computation:
\begin{lem}
\source{universal-secant-bundles-syzygies-canonical-curves}
\notready \label{det-gamma}
We have $\det \Gamma \simeq {q'}^* \mathcal{O}_{\PP}(k-1)(-D-{p'}^*\Delta)$. In particular, we have a natural isomorphism $\mathrm{H}^2(B, \wedge^{k+2} \Gamma ({p'}^*\Delta)(D))\simeq \mathrm{Sym}^{k-1}\mathrm{H}^0(X,E)^{\vee}$.
\end{lem}
\begin{proof}
We have $\det \Gamma \simeq {p'}^*(L^{\vee}) \otimes \det S^{\vee}$ so $\det \Gamma \simeq {q'}^* \mathcal{O}_{\PP}(k-1)(-D-{p'}^*\Delta)$ and $\wedge^{k+2} \Gamma ({p'}^*\Delta)\simeq {q'}^* \mathcal{O}_{\PP}(k-1)(-D)$. Thus $\mathrm{H}^2(B, \wedge^{k+2} \Gamma ({p'}^*\Delta)(D))\simeq \mathrm{H}^2({q'}^* \mathcal{O}_{\PP}(k-1))\simeq \mathrm{H}^2(B,{q'}^* \mathcal{O}_{\PP}(k-1)) \simeq \mathrm{Sym}^{k-1}\mathrm{H}^0(X,E)^{\vee}$.
\end{proof}


We have natural isomorphisms $$\mathrm{K}_{k-1,1}(X,-\Delta, L)\simeq \mathrm{H}^1(\wedge^k M_L(-\Delta))\simeq \mathrm{H}^0(\wedge^{k-1}M_L (L-\Delta)),$$
see \cite{lazarsfeld-VBT} or \cite[\S 3]{ein-lazarsfeld-asymptotic}. We have $\mathrm{H}^0(\wedge^{k-1}M_L (L-\Delta))^{\vee} \simeq \mathrm{H}^2(\wedge^{k+2}M_L(\Delta))$, since $\wedge^{k-1}M_L^{\vee} \simeq \wedge^{k+2}M_L(L)$.

%Further, $\mathrm{H}^2(B, \wedge^{k+2}\pi^*\mathcal{M}({p'}^*\Delta)(D)) \simeq \mathrm{H}^2(B, \wedge^{k+2}\pi^*\mathcal{M}({p'}^*\Delta)) \simeq \mathrm{H}^2(X,\wedge^{k+2}M_L(\Delta)).$

Taking exterior powers of $0 \to S \to \pi^* \mathcal{M} \to \Gamma \to 0$ and twisting by $\mathcal{O}_{B}(D+{p'}^*\Delta)$
induces
$$\phi \; : \; \mathrm{H}^2(\wedge^{k+2}M_L(\Delta)) \to \mathrm{H}^2(\wedge^{k+2} \Gamma ({p'}^*\Delta)(D))\simeq \mathrm{Sym}^{k-1}\mathrm{H}^0(E)^{\vee}.$$
We will show that $\phi$ is surjective. We write $\mathcal{O}_B(j)$ for ${q'}^*\mathcal{O}_{\PP}(j) $. Consider the exact sequence
{\small{$$\ldots \to \bigwedge^{k+1}\pi^* \mathcal{M} \otimes S({p'}^*\Delta)(D)\to  \bigwedge^{k+2}\pi^* \mathcal{M} ({p'}^*\Delta)(D) \to \bigwedge^{k+2} \Gamma ({p'}^*\Delta)(D) \to 0.$$}}
The following proposition should be compared to \cite[Lemma 6]{V2}.
\begin{prop}
\source{universal-secant-bundles-syzygies-canonical-curves}
\notready \label{most-of-van}
We have $\mathrm{H}^{2+i}(\bigwedge^{k+2-i} \pi^* \mathcal{M} \otimes \mathrm{Sym}^i(S) ({p'}^*\Delta)(D))=0$ for $1 \leq i \leq k+2$, $i \neq k, k+1$.
\end{prop}
\begin{proof}
The Euler sequence $0 \to \mathcal{O}_B(-2) \to \chi^* \mathcal{H} (-1)  \to S \to 0$
gives exact sequences
$$0 \to \chi^* \mathrm{Sym}^{i-1} \mathcal{H} (-i-1) \to \chi^*\mathrm{Sym}^{i} \mathcal{H} (-i) \to \mathrm{Sym}^i(S) \to 0.$$
 We first claim
  $$\mathrm{H}^{2+i}(\bigwedge^{k+2-i} \pi^* \mathcal{M} \otimes \chi^*\mathrm{Sym}^{i} \mathcal{H} (-i)({p'}^*\Delta)(D))=0, \; \; \text{for $1 \leq i \leq k-1$}.$$
 The fibres of $\chi$ are projective spaces of dimension $k-1$ and $D$ has degree one on these fibres. Hence $R^j \chi_*\mathcal{O}_B(-i)(D)=0$ for all $j$ and $2 \leq i \leq k-1$. So the claim holds for $2 \leq i \leq k-1$ by the Leray spectral sequence.

 For $i=1$, the claim states $\mathrm{H}^3(\bigwedge^{k+1} \pi^* \mathcal{M} \otimes \chi^*\mathcal{H} (-1)({p'}^*\Delta)(D))=0$. We have $\omega_B \simeq {q'}^* \omega_{\PP}(3D)$, so that this is equivalent to
$$\mathrm{H}^{k+1}(\bigwedge^{k+1} \pi^* \mathcal{M}^{\vee} \otimes \chi^*\mathcal{H}^{\vee}\otimes {q'}^*(\omega_{\PP}(1)) (-{p'}^*\Delta)(2D))=0$$
The claim now follows from the exact sequence $ 0 \to \mathcal{O}_{B}(1)(-D) \to \chi^* \mathcal{H}^{\vee} \to {p'}^*N^{\vee} \to 0$ since
{\small{\begin{align*}
&\mathrm{H}^{k+1}(B, \bigwedge^{k+1} \pi^* \mathcal{M}^{\vee} \otimes {q'}^*(\omega_{\PP}(2)) (-{p'}^*\Delta)(D)) \simeq \mathrm{H}^{k+1}(X\times \PP, \bigwedge^{k+1} M_L^{\vee}(-\Delta)\boxtimes \omega_{\PP}(2))=0, \; \; \text{and} \\
&\mathrm{H}^{k+1}(B,\bigwedge^{k+1} \pi^* \mathcal{M}^{\vee} \otimes {p'}^*N^{\vee} \otimes {q'}^*(\omega_{\PP}(1)) (-{p'}^*\Delta)(2D))  \simeq \mathrm{H}^{k+1}(X \times \PP, \bigwedge^{k+1} M_L^{\vee}\otimes N^{\vee}(-\Delta) \boxtimes \omega_{\PP}(1)),
\end{align*}}}
by the K\"unneth formula.


To finish the proof in the case $1 \leq i \leq k-1$,  it suffices to note that
 $$\mathrm{H}^{3+i}(\bigwedge^{k+2-i} \pi^* \mathcal{M} \otimes \chi^*\mathrm{Sym}^{i-1} \mathcal{H} (-i-1)({p'}^*\Delta)(D))=0, \; \; \text{for $1 \leq i \leq k-1$},$$ is immediate as above. We are now left with the case $i=k+2$. The inclusion $$\pi^*( \mathrm{Sym}^{k+2}N(\Delta) \boxtimes \mathcal{O}_{\PP}(-(k+2))) \hookrightarrow \mathrm{Sym}^{k+2}(S)({p'}^*\Delta)$$ is an isomorphism off $D$. Since $\dim D=k+3$, it suffices to show $$\mathrm{H}^{k+4}(X \times \PP^{k+2}, \mathrm{Sym}^{k+2}N(\Delta) \boxtimes \mathcal{O}_{\PP}(-(k+2)))=0.$$ This follows from the K\"unneth formula. \end{proof}

We next show that we continue to have the above vanishing in the case $i=k+1$. The following lemma, stated in \cite{V1}, Proof of Prop.\ 6, was explained to us by C.\ Voisin.
\begin{lem}
\source{universal-secant-bundles-syzygies-canonical-curves}
\notready \label{LM-mult-map}
The multiplication map $\mathrm{H}^0(X,E(-\Delta)) \otimes \mathrm{H}^0(X,L-\Delta) \to \mathrm{H}^0(X,E(L-2\Delta))$ is surjective.
\end{lem}
\begin{proof}
Let $C \in |L-\Delta|$ be a smooth curve. It suffices to prove surjectivity of the restricted multiplication map $\mathrm{H}^0(E_{|_C}(-\Delta)) \otimes \mathrm{H}^0(\mathrm{K}_C) \to \mathrm{H}^0(E_{|_C}(\mathrm{K}_C-\Delta))$, \cite[Observation 2.3]{gallego-purnaprajna}. Now $\mathrm{H}^0(E_{|_C}(-\Delta)) \simeq \mathrm{H}^0(B) \oplus \mathrm{H}^0(\mathrm{K}_C-B)$ where $B$ is a $g^1_k$ on $C$, \cite{voisin-wahl}. The statement now follows from the following: for any base-point free line bundle $A$ on $C$ with $h^0(A) \geq 2$, the multiplication map $\mathrm{H}^0(\mathrm{K}_C) \otimes \mathrm{H}^0(A) \to \mathrm{H}^0(\mathrm{K}_C+A)$ is surjective. If $h^0(A)=2$, this follows immediately from the base-point free pencil trick. Indeed, we have an exact sequence
$$0 \to K_C \otimes A^{-1} \to H^0(A) \otimes K_C \to K_C \otimes A \to 0,$$
and the claim follows from fact that $h^1(K_C \otimes A^{-1})=h^1(H^0(A) \otimes K_C )=h^0(A)$. Otherwise, let $Z$ be a general effective divisor of degree $h^0(A)-2$. Thus $\mathrm{H}^0(\mathrm{K}_C) \otimes \mathrm{H}^0(A-Z) \to \mathrm{H}^0(\mathrm{K}_C+A-Z)$ is surjective. Since this holds for any general such divisor, this proves the claim.
\end{proof}
We will make use of the following direct consequence of the previous lemma.
\begin{lem}
\source{universal-secant-bundles-syzygies-canonical-curves}
\notready \label{mult-maps}
We have $\mathrm{H}^2(X,M_L(\Delta) \otimes N)=\mathrm{H}^2(M_L(\Delta))=0$.
\end{lem}
\begin{proof}
From the exact sequence
$$0 \to N \to \mathrm{H}^0(E(-\Delta)) \otimes \mathcal{O}_X \to E(-\Delta)\to 0,$$
it suffices to show $\mathrm{H}^2(M_L(\Delta))=0$ and $\mathrm{H}^1(M_L(E))=0$. The first vanishing follows immediately from the defining sequence for $M_L$.

The vanishing $\mathrm{H}^1(M_L(E))=0$ is equivalent to surjectivity of the multiplication map $\mathrm{H}^0(L) \otimes \mathrm{H}^0(E) \to \mathrm{H}^0(E(L))$. We will use the following principle. If we have a commutative diagram
{\small{$$\text{[Diagram omitted; see source PDF.]}$$}}
of vector spaces with exact rows, and if $h$ is surjective, then so is $\delta$. If, further, $f$ is surjective, then so is $g$.

We have surjectivity of $\mathrm{H}^0(L-\Delta) \otimes \mathrm{H}^0(E(-\Delta)) \to \mathrm{H}^0(E(L-2\Delta))$ by Lemma \ref{LM-mult-map}.  For a \emph{general} $s \in \mathrm{H}^0(E)$, we have the exact sequence $0 \to \mathcal{O}_X \to E \to L'\otimes I_{Z(s)} \to 0$. This implies $E_{|_{\Delta}} \simeq \mathcal{O}_{\Delta}^{\oplus 2}$ is trivial. Thus $\mathrm{H}^0(L-\Delta) \otimes \mathrm{H}^0(E_{|_{\Delta}}) \twoheadrightarrow \mathrm{H}^0(E_{|_{\Delta}}(L-\Delta))$. This shows $\mathrm{H}^0(L-\Delta) \otimes \mathrm{H}^0(E) \twoheadrightarrow \mathrm{H}^0(E(L-\Delta))$. Since $\mathrm{H}^0(L_{|_{\Delta}}) \otimes \mathrm{H}^0(E) \twoheadrightarrow \mathrm{H}^0(E_{|_{\Delta}}(L))$, using the triviality of $E_{|_{\Delta}}$, we obtain the vanishing $\mathrm{H}^1(M_L(E))=0$.\end{proof}

We now prove the vanishing in case $i=k+1$. The following proof should be compared to \cite[Lemma 8]{V2}.
\begin{prop}
\source{universal-secant-bundles-syzygies-canonical-curves}
\notready \label{green-odd-main}
We have $\mathrm{H}^{k+3}(B,\pi^* \mathcal{M} \otimes \mathrm{Sym}^{k+1}(S) ({p'}^*\Delta)(D))=0$.
 \end{prop}
\begin{proof}
We have the short exact sequence
$$0 \to \chi^* \mathrm{Sym}^{k} \mathcal{H} (-k-2) \to \chi^*\mathrm{Sym}^{k+1} \mathcal{H} (-k-1) \to \mathrm{Sym}^{k+1}(S) \to 0.$$
We firstly claim that $$\mathrm{H}^{k+3}(\pi^* \mathcal{M} \otimes  \chi^* \mathrm{Sym}^{k} \mathcal{H} (-k-2)  ({p'}^*\Delta)(D))\to \mathrm{H}^{k+3}(\pi^* \mathcal{M} \otimes  \chi^*\mathrm{Sym}^{k+1} \mathcal{H} (-k-1)  ({p'}^*\Delta)(D))$$
is surjective. We have $\mathcal{O}(D) \simeq \mathcal{O}(1)\otimes \chi^*\mathcal{O}_{\PP(\mathcal{F})}(-1)$ and further $\omega_{\chi}\simeq \chi^*\det \mathcal{H}^{\vee}(-k)$. The map can thus be written as
{\small{\begin{align*}
&\mathrm{H}^4(\PP(\mathcal{F}),f^*M_L(\Delta) \otimes \mathrm{Sym}^k\mathcal{H} \otimes R^{k-1}\chi_*\mathcal{O}(-(k+1))\otimes \mathcal{O}_{\PP(\mathcal{F})}(-1)) \to  \\
&\mathrm{H}^4(\PP(\mathcal{F}),f^*M_L(\Delta) \otimes \mathrm{Sym}^{k+1}\mathcal{H} \otimes R^{k-1}\chi_*\mathcal{O}(-k)\otimes \mathcal{O}_{\PP(\mathcal{F})}(-1)) .
\end{align*}}}
Using relative duality, this becomes
{\small{\begin{align*}
&\mathrm{H}^4(\PP(\mathcal{F}),f^*M_L(\Delta) \otimes \mathrm{Sym}^k\mathcal{H} \otimes \mathcal{H} \otimes \mathcal{O}_{\PP(\mathcal{F})}(-1)\otimes \det \mathcal{H}) \to  \\
&\mathrm{H}^4(\PP(\mathcal{F}),f^*M_L(\Delta) \otimes \mathrm{Sym}^{k+1}\mathcal{H} \otimes \mathcal{O}_{\PP(\mathcal{F})}(-1)\otimes \det \mathcal{H})),
\end{align*}}}
since $\chi_*\mathcal{O}(n)\simeq \mathrm{Sym}^n\mathcal{H}^{\vee}$. This map is surjective, since the composite $ \mathrm{Sym}^{k+1}\mathcal{H} \to \mathrm{Sym}^k\mathcal{H} \otimes \mathcal{H} \to \mathrm{Sym}^{k+1}\mathcal{H} $ of natural maps is multiplication by $k+1$.

To conclude, it suffices that $\mathrm{H}^{k+4}(\pi^* \mathcal{M} \otimes  \chi^* \mathrm{Sym}^{k} \mathcal{H} (-k-2)({p'}^*\Delta)(D))=0$, or, equivalently
$$\mathrm{H}^5(\PP(\mathcal{F}),f^*M_L(\Delta) \otimes \mathrm{Sym}^k\mathcal{H} \otimes \mathcal{H} \otimes \mathcal{O}_{\PP(\mathcal{F})}(-1)\otimes \det \mathcal{H}))=0.$$
The same argument identifies this space with $\mathrm{H}^{k+4}(\pi^* \mathcal{M} \otimes \chi^*\mathcal{H} (-2k-1+{p'}^*\Delta)(D))$, using $\chi_*\mathcal{O}(n)\simeq \mathrm{Sym}^n\mathcal{H}^{\vee}$ again. From the exact sequence
$$0 \to {p'}^*N \to \chi^* \mathcal{H} \to {q'}^*\mathcal{O}_{\PP}(-1)(D) \to 0,$$
the vanishing is implied by
{\small{\begin{align*}
\mathrm{H}^{k+4}(X \times \PP, M_L(\Delta) \otimes N \boxtimes \mathcal{O}_{\PP}(-2k-1))=\mathrm{H}^{k+4}(X \times \PP, M_L(\Delta)\boxtimes \mathcal{O}_{\PP}(-2k-2))=0
\end{align*}}}
which follow from Lemma \ref{mult-maps} and the K\"unneth formula.
\end{proof}

From the above vanishings, to show the surjectivity of $\phi$ it is now enough to show that the natural map $\mathrm{H}^{k+2}(\pi^*\mathcal{M} \otimes \mathrm{Sym}^{k+1}(S)({p'}^*\Delta)(D)) \to \mathrm{H}^{k+2}(\bigwedge^2 \pi^*\mathcal{M} \otimes \mathrm{Sym}^{k}(S)({p'}^*\Delta)(D))$ is surjective. We first rewrite this map.
\begin{lem}
\source{universal-secant-bundles-syzygies-canonical-curves}
\notready \label{identilem}
For any $j$, we have have natural isomorphisms
{\small{
\begin{align*}
\mathrm{H}^j(B, \pi^* \mathcal{M} \otimes \mathrm{Sym}^{k+1}(S)({p'}^*\Delta)(D))&\simeq \mathrm{H}^{j+1}(B, \pi^* \mathcal{M} \otimes S(-2k)({p'}^*\Delta)(D)), \\
\mathrm{H}^{j}(B,\bigwedge^{2}\pi^*\mathcal{M}  \otimes \mathrm{Sym}^k(S)({p'}^*\Delta)(D))&\simeq \mathrm{H}^{j+1}(B,\bigwedge^{2}\pi^*\mathcal{M}(-2k)({p'}^*\Delta)(D))
\end{align*}
}}
\end{lem}
\begin{proof}
Arguing as in Proposition \ref{green-odd-main}, identify $\mathrm{H}^j(B, \pi^* \mathcal{M} \otimes \mathrm{Sym}^{k+1}(S)({p'}^*\Delta)(D))$ with
{\small{\begin{align*}
\mathrm{Ker}(&\mathrm{H}^{j+2-k}(\PP(\mathcal{F}),f^*M_L(\Delta) \otimes \mathrm{Sym}^k\mathcal{H} \otimes \mathcal{H} \otimes \mathcal{O}_{\PP(\mathcal{F})}(-1) \otimes \det\mathcal{H}) \twoheadrightarrow  \\
&\mathrm{H}^{j+2-k}(\PP(\mathcal{F}),f^*M_L(\Delta) \otimes \mathrm{Sym}^{k+1}\mathcal{H} \otimes \mathcal{O}_{\PP(\mathcal{F})}(-1)) \otimes \det \mathcal{H}) .
\end{align*}}}
The above surjection splits naturally, identifying $\mathrm{H}^j(B, \pi^* \mathcal{M} \otimes \mathrm{Sym}^{k+1}(S)({p'}^*\Delta)(D))$ with
{\small{\begin{align*}
\mathrm{Coker}(&\mathrm{H}^{j+2-k}(\PP(\mathcal{F}),f^*M_L(\Delta) \otimes \mathrm{Sym}^{k+1}\mathcal{H} \otimes \mathcal{O}_{\PP(\mathcal{F})}(-1) \otimes \det \mathcal{H}) \hookrightarrow \\
&\mathrm{H}^{j+2-k}(\PP(\mathcal{F}),f^*M_L(\Delta) \otimes \mathrm{Sym}^k\mathcal{H} \otimes \mathcal{H} \otimes \mathcal{O}_{\PP(\mathcal{F})}(-1)) \otimes \det \mathcal{H}),
\end{align*}}}
which is naturally identified with
\begin{align*}
\mathrm{Coker}\left( \mathrm{H}^{j+1}(B, \pi^* \mathcal{M}(-2k-2)({p'}^*\Delta)(D)) \hookrightarrow
 \mathrm{H}^{j+1}(B, \pi^* \mathcal{M} \otimes \mathcal{H}(-2k-1)({p'}^*\Delta)(D))\right).
\end{align*}
From the sequence $0 \to \mathcal{O}_B(-2) \to \chi^* \mathcal{H} (-1)  \to S \to 0$, the above space is isomorphic to $\mathrm{H}^{j+1}(B, \pi^* \mathcal{M} \otimes S(-2k)({p'}^*\Delta)(D))$, as required.

For the second isomorphism, the exact sequence
$$0 \to \chi^* \mathrm{Sym}^{k-1} \mathcal{H} (-k-1) \to \chi^*\mathrm{Sym}^{k} \mathcal{H} (-k) \to \mathrm{Sym}^k(S) \to 0,$$
together with relative duality for $\chi$, provides natural isomorphisms
{\small{\begin{align*}
\mathrm{H}^{j}(B, \bigwedge^{2}\pi^*\mathcal{M}  \otimes \mathrm{Sym}^k(S)({p'}^*\Delta)(D))& \simeq \mathrm{H}^{j+1}(B,\bigwedge^{2}\pi^*\mathcal{M}  \otimes \chi^*\mathrm{Sym}^{k-1}(\mathcal{H})(-k-1)({p'}^*\Delta)(D)) \\
&\simeq \mathrm{H}^{k+1-j}(\PP(\mathcal{F}), \bigwedge^{2}f^*M_L(\Delta) \otimes \mathrm{Sym}^{k-1}\mathcal{H} \otimes \mathcal{O}_{\PP(\mathcal{F})}(-1)\otimes \det \mathcal{H})\\
& \simeq \mathrm{H}^{j+1}(B,\bigwedge^{2}\pi^*\mathcal{M}  \otimes {q'}^*\mathcal{O}_{\PP}(-2k)({p'}^*\Delta)(D)).
\end{align*}}}
\end{proof}
In order to prove surjectivity of the natural map $$\mathrm{H}^{k+3}(\pi^* \mathcal{M} \otimes S(-2k)({p'}^*\Delta)(D)) \to \mathrm{H}^{k+3}(\bigwedge^{2}\pi^*\mathcal{M}(-2k)({p'}^*\Delta)(D)),$$ we follow a strategy reminiscent of \cite[\S 3, Fourth step]{V2}. Consider the exact sequence
$$0 \to \mathrm{Sym}^2(S) \to \pi^*\mathcal{M} \otimes S \to \bigwedge^2 \pi^*\mathcal{M} \to \bigwedge^2 \Gamma \to 0.$$
We have $\mathrm{H}^{k+4}(\pi^* \mathcal{M} \otimes S(-2k)({p'}^*\Delta)(D))=0$ by Lemma \ref{identilem} and Proposition \ref{green-odd-main}. Thus the surjectivity of $\phi$ will follow from
$$\mathrm{H}^{k+4}(\mathrm{Sym}^2(S)(-2k)({p'}^*\Delta)(D))=\mathrm{H}^{k+3}(\bigwedge^2\Gamma (-2k)({p'}^*\Delta)(D))=0.$$
\begin{lem}
\source{universal-secant-bundles-syzygies-canonical-curves}
\notready
We have $\mathrm{H}^{k+4}(\mathrm{Sym}^2(S)(-2k)({p'}^*\Delta)(D))=0$.
\end{lem}
\begin{proof}
The inclusion $\pi^*( \mathrm{Sym}^{2}N(\Delta) \boxtimes \mathcal{O}_{\PP}(-2)) \hookrightarrow \mathrm{Sym}^{2}(S)({p'}^*\Delta)$ is an isomorphism off $D$, so that it suffices to show $\mathrm{H}^{k+4}(X \times \PP, \mathrm{Sym}^{2}N(\Delta) \boxtimes \mathcal{O}_{\PP}(-2-2k))=0$. By the K\"unneth formula, it in turn suffices to show $\mathrm{H}^2(X,\mathrm{Sym}^{2}N(\Delta))=0$. We have an exact sequence
$$0 \to  \mathrm{Sym}^2N \to \mathrm{Sym}^2 \mathrm{H}^0(E(-\Delta)) \otimes \mathcal{O}_X \to \mathrm{H}^0(E(-\Delta)) \otimes E(-\Delta) \to L-\Delta \to 0.$$
Since $\mathrm{H}^2(X,\mathcal{O}_X(\Delta))=H^1(X,E)=0$, it is enough to show that the determinant map
$$\mathrm{det} : \mathrm{H}^0(E(-\Delta)) \otimes \mathrm{H}^0(E) \to \mathrm{H}^0(L)$$
is surjective. Let $C \in |L+\Delta|$ be general (in particular $C$ is disjoint from $\Delta$) and let $A$ be any $g^1_{k+2}$ on $C$. We have an exact sequence
$$0 \to \mathrm{H}^0(A)^{\vee} \otimes \mathcal{O}_X \to E \to \omega_C \otimes A^{\vee} \to 0,$$
\cite{lazarsfeld-BNP}. Restricting to $C$ we have isomorphisms $\mathrm{H}^0(E)=\mathrm{H}^0(E_{|_C})\simeq \mathrm{H}^0(C,A)\oplus \mathrm{H}^0(\omega_C \otimes A^{\vee})$, \cite{voisin-wahl}. Twisting the above exact sequence by $-\Delta$, we have a natural surjection $\mathrm{H}^0(E(-\Delta)) \twoheadrightarrow \mathrm{H}^0(\omega_C \otimes A^{\vee})$. Further, restriction provides an isomorphism $\mathrm{H}^0(X,L) \simeq \mathrm{H}^0(C,\omega_C)$. Thus it suffices to show that the Petri map
$$\mathrm{H}^0(A) \otimes \mathrm{H}^0(\omega_C \otimes A^{\vee}) \to \mathrm{H}^0(\omega_C),$$
is surjective. By Lemma \ref{BNP-gen}, $C$ is Petri general, so the Petri map is injective. But both sides have the same dimension, by the Riemann--Roch theorem, so the Petri map is an isomorphism.

\end{proof}
We now complete the proof of the surjectivity of $\phi$.
\begin{prop}
\source{universal-secant-bundles-syzygies-canonical-curves}
\notready
We have $\mathrm{H}^{k+3}(\bigwedge^2\Gamma (-2k)({p'}^*\Delta)(D))=0$. In particular, $\phi$ is surjective.
\end{prop}
\begin{proof}
We use the exact sequence $0 \to \mathcal{O}_B \to \mathcal{O}_B(D) \to \mathcal{O}_D(D) \to 0.$ We first show \begin{equation} \label{rest-D}
\source{universal-secant-bundles-syzygies-canonical-curves}
\mathrm{H}^{k+2}(D,\bigwedge^2\Gamma (-2k)({p'}^*\Delta)(2D))=0.
\end{equation}
We have the exact sequence
$$0 \to \mathrm{Sym}^2(S)_{|_D} \to \pi^*\mathcal{M} \otimes S_{|_D} \to \bigwedge^2 \pi^*\mathcal{M}_{|_D} \to \bigwedge^2 \Gamma_{|_D} \to 0.$$
Since $D \to \mathcal{Z}$ is a projective bundle with three dimensional fibres, and $\mathcal{O}_D(D)$ has degree $-1$ on the fibres, $\mathrm{H}^j(D,\pi^*A(nD))=0$
for any bundle $A$ on $X \times \PP$, any integer $j$ and $1 \leq n \leq 3$. Thus $\mathrm{H}^{k+2}(D,\bigwedge^2\pi^*\mathcal{M}(-2k)({p'}^*\Delta)(2D))=0.$ Next, we have $\mathrm{H}^{k+3}(D,\pi^*\mathcal{M} \otimes S(-2k)({p'}^*\Delta)(2D))\simeq \mathrm{H}^{k+3}(D,\pi^*\mathcal{M} \otimes \chi^*\mathcal{H}(-2k-1)({p'}^*\Delta)(2D))$ from the exact sequence $$0 \to \mathcal{O}_D(-2) \to \chi^*\mathcal{H}_{|_D}(-1) \to S_{|_D} \to 0.$$ This space vanishes from the exact sequence $0 \to {p'}^*N \to \chi^*\mathcal{H} \to \mathcal{O}_B(-1)(D) \to 0.$ Lastly, $\mathrm{H}^{k+4}(D, \mathrm{Sym}^2(S)(-2k)({p'}^*\Delta)(2D))=0$ since $\dim D=k+3$, giving the vanishing (\ref{rest-D}).

To finish the proof, it suffices to show $\mathrm{H}^{k+3}(\bigwedge^2\Gamma (-2k)({p'}^*\Delta)(2D))=0$. We have $$\bigwedge^2 \Gamma \simeq \bigwedge^k \Gamma^{\vee}(\det \Gamma)\simeq \bigwedge^k \Gamma^{\vee}(k-1)(-{p'}^*\Delta-D)$$ by Lemma \ref{det-gamma}. Thus, $\mathrm{H}^{k+3}(\bigwedge^2\Gamma (-2k)({p'}^*\Delta)(2D))$ is Serre dual to $\mathrm{H}^{1}(\bigwedge^k\Gamma (2D-2))$. From the exact sequence
{\small{$$0 \to \mathrm{Sym}^k S \to \ldots \to \bigwedge^{k-1}\pi^*\mathcal{M} \otimes S \to \bigwedge^k \pi^* \mathcal{M} \to \bigwedge^k \Gamma \to 0,$$}}
it suffices to show $\mathrm{H}^{1+i}(\bigwedge^{k-i}\pi^*\mathcal{M} \otimes \mathrm{Sym}^i(S)(2D-2))=0$ for $0 \leq i \leq k$. In the case $i=0$ we have $\mathrm{H}^{1}(\bigwedge^{k}\pi^*\mathcal{M}(2D-2))=0$ immediately from the K\"unneth formula on $X \times \PP$. For $1 \leq i \leq k-2$ the required vanishing follows immediately from the exact sequence
$$0 \to \chi^* \mathrm{Sym}^{i-1} \mathcal{H} (-i-1) \to \chi^*\mathrm{Sym}^{i} \mathcal{H} (-i) \to \mathrm{Sym}^i(S) \to 0,$$
as in the proof of Proposition \ref{most-of-van}.

It remains to deal with the cases $i=k-1,k$. For $i=k-1$, we have
{\small{
\begin{align*}
\mathrm{H}^{k}(B,\pi^*\mathcal{M} \otimes \mathrm{Sym}^{k-1}(S)(2D-2)) & \simeq \mathrm{H}^{k+1}(B,\pi^*\mathcal{M} \otimes \chi^*\mathrm{Sym}^{k-2}(\mathcal{H})(-k)(2D-2))\\
&\simeq \mathrm{H}^2(\PP(\mathcal{F}),f^*M_L \otimes \mathrm{Sym}^{k-2}(\mathcal{H})\otimes R^{k-1}\chi_*\mathcal{O}_B(-k)\otimes\mathcal{O}_{\PP(\mathcal{F})}(-2))\\
&\simeq \mathrm{H}^2(\PP(\mathcal{F}),f^*M_L \otimes \mathrm{Sym}^{k-2}(\mathcal{H})\otimes\det \mathcal{H}\otimes\mathcal{O}_{\PP(\mathcal{F})}(-2)),
\end{align*}
}}
where we used $\omega_{\chi} \simeq \chi^*\det \mathcal{H}^{\vee}(-k)$ and relative duality for the last equality. Since $\mathrm{Sym}^{k-2}\mathcal{H} \simeq (\chi_*\mathcal{O}_B(k-2))^{\vee} \simeq R^{k-1}\chi_*\mathcal{O}_B(2-2k)\otimes \det \mathcal{H}^{\vee}$, we may identify the above space with $$\mathrm{H}^{k+1}(B, \pi^*\mathcal{M}(-2k)(2D))=0,$$ by the K\"unneth formula.

To complete the proof, we need to show $\mathrm{H}^{k+1}(\mathrm{Sym}^k(S)(2D-2))=0$. As in the proof of Proposition \ref{green-odd-main}, the map
$$\mathrm{H}^{k+1}(\mathrm{Sym}^{k-1}(\mathcal{H})(-k-1)(2D-2))\to \mathrm{H}^{k+1}(\mathrm{Sym}^{k}(\mathcal{H})(-k)(2D-2))$$
is surjective, so $\mathrm{H}^{k+1}(\mathrm{Sym}^k(S)(2D-2))$ is isomorphic to
$$ \mathrm{Ker}\left(  \mathrm{H}^{k+2}(B,\chi^*\mathrm{Sym}^{k-1}\mathcal{H}(-k-1)(2D-2))\twoheadrightarrow \mathrm{H}^{k+2}(B,\chi^*\mathrm{Sym}^{k}\mathcal{H}(-k)(2D-2) \right)$$
which is naturally identified with
{\small{
\begin{align*} \mathrm{Ker}(  &\mathrm{H}^{3}(\PP(\mathcal{F}),\mathrm{Sym}^{k-1}\mathcal{H}\otimes R^{k-1}\chi_*\mathcal{O}_B(-k-1)\otimes \mathcal{O}_{\PP(\mathcal{F})}(-2))\twoheadrightarrow \\
& \mathrm{H}^{3}(\PP(\mathcal{F}), \mathrm{Sym}^{k}\mathcal{H} \otimes R^{k-1}\chi_*\mathcal{O}_B(-k)\otimes \mathcal{O}_{\PP(\mathcal{F})}(-2) ).
\end{align*}
}}
As in the proof of Lemma \ref{identilem}, the above surjection splits and, using relative duality, the above space is identified with
$$\mathrm{Coker}\left( \mathrm{H}^{k+2}(B,\mathcal{O}_B(-2k-2)(2D)) \to \mathrm{H}^{k+2}(B,\mathcal{O}_B(-2k-1)\otimes \chi^* \mathcal{H}(2D))\right).$$
From the sequence $0 \to {p'}^*N \to \chi^*\mathcal{H} \to \mathcal{O}_B(-1)(D) \to 0$, and since $\mathrm{H}^0(X,N)=\mathrm{H}^1(X,N)=0$, we have $\mathrm{H}^{k+2}(\mathcal{O}_B(-2k-1)\otimes \chi^* \mathcal{H}(2D)) \simeq \mathrm{H}^{k+2}(\mathcal{O}_B(-2k-2)(3D))$. The map
$ \mathrm{H}^{k+2}(\mathcal{O}_B(-2k-2)(2D)) \to  \mathrm{H}^{k+2}(\mathcal{O}_B(-2k-2)(3D))$ is multiplication by the section $s \in \mathrm{H}^0(\mathcal{O}_B(D))$ induced from the composition $\mathcal{O}_B \hookrightarrow \chi^*\mathcal{H}(1)$, coming from the Euler sequence, with the map $\chi^*\mathcal{H}(1) \twoheadrightarrow \mathcal{O}_B(D)$ coming from the definition of $\mathcal{H}$. Note $s \neq 0$ as $\mathrm{H}^1(B,{p'}^*N(1))=0$, so this map is an isomorphism (by Equation (\ref{mult-D-iso})). This completes the proof.



\end{proof}
\begin{rem}
\source{universal-secant-bundles-syzygies-canonical-curves}
\notready
Note that we do not need the geometry of Grassmannians, unlike \cite[\S 3, Fourth step]{V2}.
\end{rem}

Let $\psi^{\vee}: \mathrm{Sym}^{k-1}\mathrm{H}^0(E) \xrightarrow{\sim} \mathrm{K}_{k,1}(X,L+\Delta)$ and $\phi^{\vee}: \mathrm{Sym}^{k-1}\mathrm{H}^0(E) \hookrightarrow \mathrm{K}_{k-1,1}(X,-\Delta, L)$ be the duals to the maps from Lemma \ref{green-pic2} and Proposition \ref{green-odd-main}. By Lemma \ref{proj-crit}, to complete the proof of Voisin's Theorem, it only remains to show that $\phi^{\vee}=pr_k \circ \psi^{\vee}$.

By duality and from the sequence $0 \to \mathcal{S} \to \pi^* \mathcal{M} \to \Gamma \to 0,$ we identify $\phi^{\vee}$ with
$$\mathrm{H}^{k+2}(B, \wedge^{k-1} \mathcal{S}\otimes \pi^*(L-\Delta \boxtimes \omega_{\PP})(2D)) \to \mathrm{H}^{k+2}(B, \pi^*\wedge^{k-1}M_L(L-\Delta) \boxtimes \omega_{\PP})(2D)),$$
which can, in turn, be identified with the natural map
$$\phi^{\vee} \; : \; \mathrm{H}^{k+2}(X \times \PP, \wedge^{k-1}N (L-\Delta) \boxtimes \omega_{\PP}(1-k)) \to \mathrm{H}^{k+2}(X \times \PP, \wedge^{k-1}M_L(L-\Delta) \boxtimes \omega_{\PP}).$$
\begin{thm}
\source{universal-secant-bundles-syzygies-canonical-curves}
\notready
With notation as above, we have $\mathrm{K}_{k,1}(X,L)=0$.
\end{thm}
\begin{proof}
We need to show $\phi^{\vee}=pr_k \circ \psi^{\vee}$. To relate $\psi^{\vee}$ and $\phi^{\vee}$, let $\tilde{\pi}: \widetilde{B} \to X \times \PP$ denote the blow-up in the codimension two locus $\widetilde{\mathcal{Z}}:=\{ (x,s) \; | s(x)=0 \}$. Let $\widetilde{D}$ denote the exceptional divisor. Define $\tilde{p}:=p \circ \tilde{\pi}$ and $\tilde{q}:=q \circ \tilde{\pi}$. We have the exact sequence
$$0 \to \tilde{\pi}^*(\mathcal{O}_X(-\Delta) \boxtimes \mathcal{O}_{\PP}(-2))(\widetilde{D}) \to \tilde{\pi}^*(E(-\Delta) \boxtimes \mathcal{O}_{\PP}(-1)) \to {\tilde{p}}^* L \otimes I_{\widetilde{D}} \to 0.$$
We have $\tilde{q}_*({\tilde{p}}^* L \otimes I_{\widetilde{D}}) \simeq \mathrm{H}^0(E(-\Delta)) \otimes \mathcal{O}_{\PP}(-1)$. Define $\tilde{S}_{L'}$ and $\tilde{S}_L$ by exact sequences
{\small{\begin{align*}
0 \to &\tilde{S}_{L'} \to {\tilde{q}}^*{\tilde{q}}_*({\tilde{p}}^* {L'} \otimes I_{\widetilde{D}})  \to {\tilde{p}}^* L' \otimes I_{\widetilde{D}} \to 0, \\
0 \to &\tilde{S}_L \to \mathrm{H}^0(E(-\Delta)) \otimes {\tilde{q}}^*\mathcal{O}_{\PP}(-1) \to {\tilde{p}}^* L \otimes I_{\widetilde{D}} \to 0.
\end{align*}}}
Then $\psi^{\vee}$ is the natural map $\psi^{\vee} \; : \; \mathrm{H}^{k+3}(\wedge^{k+1}\tilde{S}_{L'} \otimes \omega_{\widetilde{B}}) \to \mathrm{H}^{k+3}(\tilde{\pi}^*(\wedge^{k+1}M_{L'} \boxtimes \mathcal{O}_{\PP}) \otimes \omega_{\widetilde{B}}).$
By taking exterior powers of the exact sequence
$$0 \to \tilde{\pi}^*(N \boxtimes \mathcal{O}_{\PP}(-1)) \to \tilde{S}_L \to \tilde{\pi}^* (\mathcal{O}_X(-\Delta) \boxtimes \mathcal{O}_{\PP}(-2))(\widetilde{D}) \to 0$$
and using $\mathrm{H}^0(X,\wedge^{k-2}N(L-2\Delta))=\mathrm{H}^0(N^{\vee}(-\Delta))=\mathrm{H}^2(N(\Delta))=0$, identify $\phi^{\vee}$ with
$$\mathrm{H}^{k+2}(\wedge^{k-1} \tilde{S}_L (L-\Delta) \otimes {\tilde{q}}^* \omega_{\PP}) \to \mathrm{H}^{k+2}(\tilde{\pi}^*\wedge^{k-1}M_L (L-\Delta) \boxtimes \omega_{\PP}),$$
induced by $\tilde{S}_L \hookrightarrow \tilde{p}^* M_L$. Using exterior powers of the defining sequence for $\tilde{S}_L$, this can be further identified with
$$\phi^{\vee} \; : \; \mathrm{H}^{k+3}(\wedge^k \tilde{S}_L(-{\tilde{p}}^*\Delta) \otimes \omega_{\widetilde{B}}) \to \mathrm{H}^{k+3}(\tilde{\pi}^*(\wedge^k M_L(-\Delta) \boxtimes \mathcal{O}_{\PP}) \otimes \omega_{\widetilde{B}}),$$
using $\mathrm{H}^0(\mathcal{O}_X(-\Delta))=\mathrm{H}^1(\mathcal{O}_X(-\Delta))=0$.


Let $U \seq \widetilde{B}$ be the complement of the codimension two locus $\tilde{\pi}^{-1}(X \times \PP(\mathrm{H}^0(E(-\Delta)))$. We have an exact sequence
$0 \to \tilde{S}_{L_{|_U}} \to \tilde{S}_{L'_{|_U}} \to \mathcal{O}_U(-\tilde{p}^*\Delta) \to 0$, giving a commutative diagram
{\small{
$$\text{[Diagram omitted; see source PDF.]}$$
}}
where $q: \tilde{\pi}^*(\wedge^{k+1}M_{L'} \boxtimes \mathcal{O}_{\PP}) \to \tilde{\pi}^*(\wedge^k M_L(-\Delta) \boxtimes \mathcal{O}_{\PP})$ is the projection. This diagram extends uniquely to $\widetilde{B}$, giving the claim. \end{proof}
